El subsistema de autenticación, accesible desde la pantalla de inicio de sesión, permite a los usuarios registrarse e iniciar sesión
y obtener un token JWT para realizar las peticiones al backend. 

Por el lado del frontend, el componente LoginComponent contiene los formularios 
de registro y acceso, que se implementan mediante Angular Reactive Forms. Gracias a ello, permite emplear validadores tanto personalizados como de 
Angular, para comprobar que los campos estén rellenos correctamente. El componente invoca al servicio AuthService, para realizar la petición POST
al endpont '/users' para el registro y al endpoint '/token/login' para el acceso. Cuando se obtiene el token, se almacena 
en el almacenamiento local del navegador mediante el servicio Storage, para su uso en futuras peticiones al backend. 


Por el lado del backend, Djoser proporciona un conjunto de endpoints por defecto para la autenticación de usuarios, como resetear contraseña 
o confirmar el correo electrónico. Una vez que se accede al enpoint, Djoser valida el token JWT y, si es válido, permite el acceso a los recursos protegidos, 
en caso contrario, devuelve un error 401 Unauthorized (RNF-2).

En la figura \ref{FIG:ERROR}
se muestra un ejemplo de error al intentar acceder a un endpoint sin proporcionar el token JWT.

\begin{figure}[Prueba PostMan]{FIG:ERROR}{Solicitud GET sin token JWT}
    \image{15cm}{}{error}
\end{figure}

Mediante PostMan, se realiza un solicitud GET sin incluir ningún token JWT, lo que provoca un error 401 Unauthorized, y el mensaje de error 
del servidor.